\documentclass[12pt]{article}
\usepackage[T1]{fontenc}
\usepackage[utf8]{inputenc}
\usepackage{lmodern}
\usepackage{textcomp}
\usepackage{enumitem}
\usepackage{geometry}
\geometry{margin=1in}
\begin{document}
\begin{center}
Answer Key - Constitutional Law - Undergraduate Level Assessment 19 \
\small (Answers are ordered exactly as in the question paper.)
\end{center}

\section*{Multiple Choice}
\begin{enumerate}[label=\arabic*.]
\item \textbf{Answer:} B
\\ \textit{Options:} A) He argues than an unequal society is inevitable.; B) He claims that by giving priority to the needs of the poor, we can increase equality.; C) He asserts that we are each responsible for our poverty.; D) He rejects the idea of equality altogether.
\\ \textit{Note:} Parfit's opposition to equality is rooted in his belief that prioritizing the needs of the poor can lead to a more equitable distribution of resources, thereby enhancing overall equality rather than merely striving for equal outcomes regardless of individual circumstances.
\item \textbf{Answer:} B
\\ \textit{Options:} A) An act is jure imperii when undertaken by an international organisation; B) An act is jure imperii when undertaken in an official State capacity; C) All acts undertaken by State officials are acts jure imperii; D) An act is jure imperii when undertaken by a State corporation
\\ \textit{Note:} An act jure imperii refers to actions taken by a state or its officials that are recognized as sovereign acts, distinguishing them from private or personal actions and thereby affirming the legal immunity typically granted to states under international law.
\item \textbf{Answer:} A
\\ \textit{Options:} A) When no-one would prefer another's bundle of resources to his or her own.; B) By reference to the ownership of private property.; C) By the amount of income tax paid by individuals.; D) When the community determines that equality has been achieved.
\\ \textit{Note:} Dworkin's argument centers on the idea that equality of resources should be measured by individuals' preferences regarding their own and others' resource bundles, ensuring that no one would prefer another's resources over their own, thus achieving an equal distribution that supports overall welfare.
\item \textbf{Answer:} B
\\ \textit{Options:} A) To advance the interests of politicians.; B) To produce virtuous citizens and encourage righteousness in individuals.; C) To obtain power.; D) To secure justice by abolishing slavery.
\\ \textit{Note:} Aristotle believed that the ultimate aim of politics is to cultivate virtue and moral character among citizens, thereby promoting a just and ethical society.
\item \textbf{Answer:} C
\\ \textit{Options:} A) Armed force was prohibited; B) Armed force was permitted with no restrictions; C) Armed force was permitted subject to few restrictions; D) Armed force was not regulated under international law prior to 1945
\\ \textit{Note:} Prior to the United Nations Charter, international law allowed states to use armed force in self-defense or to resolve disputes, leading to minimal restrictions on the use of military action.
\end{enumerate}

\section*{True / False}
\begin{enumerate}[label=\arabic*.]
\setcounter{enumi}{5}
\item \textbf{Answer:} True
\\ \textit{Options:} A) True; B) False
\\ \textit{Note:} The concept of 'natural law' originated in ancient Greek philosophy, particularly through thinkers like Aristotle, who posited that certain inherent moral principles govern human behavior and can form the basis for legal systems, thereby influencing subsequent legal and philosophical thought throughout history.
\item \textbf{Answer:} False
\\ \textit{Options:} A) True; B) False
\\ \textit{Note:} Friedrich Kelsen is not considered the first jurist to study the comparative aspect of law; rather, he is primarily known for his contributions to legal positivism and the pure theory of law, while earlier scholars such as Montesquieu and Jeremy Bentham also explored comparative law concepts.
\item \textbf{Answer:} True
\\ \textit{Options:} A) True; B) False
\\ \textit{Note:} Dworkin's notion of law as integrity emphasizes that law should be interpreted based on principles of justice and fairness, which inherently rejects authoritarian interpretations that prioritize power over individual rights and moral considerations.
\item \textbf{Answer:} True
\\ \textit{Options:} A) True; B) False
\\ \textit{Note:} The Analytical School asserts that legal customs gain legitimacy and authority when they are formally acknowledged and validated through judicial decisions, thereby integrating them into the framework of law.
\item \textbf{Answer:} False
\\ \textit{Options:} A) True; B) False
\\ \textit{Note:} Codification of law can complicate legal interpretation by creating rigid frameworks that may not account for the nuances of individual cases, thereby making it less accessible to the public who may struggle to navigate the legal system.
\end{enumerate}

\section*{Long Form Response}
\begin{enumerate}[label=\arabic*.]
\setcounter{enumi}{10}
\item \textbf{Answer:} Fragmentation in international law refers to the coexistence of multiple legal regimes, such as foreign investment law and human rights law, which can develop independently and lead to divergent rules. This phenomenon poses challenges for governance, as inconsistencies between legal frameworks may create ambiguities in compliance and enforcement. For example, a foreign investment agreement may prioritize economic benefits at the expense of human rights protections, leading to conflicts in legal obligations. As these regimes operate in isolation, their lack of coordination can undermine the coherence of international law, complicating the ability of states and international bodies to navigate overlapping legal responsibilities. Ultimately, this fragmentation necessitates a reevaluation of how different legal systems interact to promote a more unified approach to international governance.
\\ \textit{Note:} correct because it accurately describes how the existence of multiple, independently developed legal regimes in international law can result in conflicting rules that complicate compliance and governance, exemplified by the tension between economic priorities and human rights protections.
\item \textbf{Answer:} Civil law is a branch of law that governs the rights and responsibilities between individuals, distinct from criminal law, which addresses offenses against the state. It establishes a framework for resolving disputes over rights, such as contracts, property, and family matters, ensuring that individuals can seek redress when their rights are violated. Mechanisms for redress in civil law include lawsuits, where an injured party can claim compensation or specific performance through the court system. This legal structure promotes social order by providing clear guidelines for behavior and recourse for those wronged, thereby reinforcing the rule of law in society. Ultimately, civil law plays a crucial role in protecting individual rights and facilitating peaceful conflict resolution.
\\ \textit{Note:} The answer correctly identifies civil law as a distinct legal framework that regulates the rights and responsibilities between individuals and outlines the processes, such as lawsuits, through which individuals can seek redress for violations of those rights.
\end{enumerate}

\vfill
\noindent\textit{End of Answer Key}
\end{document}